\chapter{Estimation Filters}
\label{chap:filters}

This chapter specifies the mathematical contract of \lupnt{}'s estimation
layer: the recursive Kalman family (\cpp{EKF}, \cpp{UKF}, \cpp{UDUEKF},
\cpp{SRIF}, \cpp{SchmidtEKF}, \cpp{UDUStochasticCloningEKF}), the batch
weighted least-squares estimator, the fixed-interval smoothers, and the
process-noise models that drive them. The recursive filters live under
\file{cpp/lupnt/numerics/filters/}; the clock process noise lives in
\file{cpp/lupnt/dynamics/clock_dynamics.cc}. The linear-algebra kernels follow
\citet{bierman1977} and \citet{thornton1980} for the UD factorization,
\citet{tapley2004} for the square-root information filter, and
\citet{julier2004ukf} for the unscented transform. The stochastic-cloning UD
filter and the delayed-state smoother reproduce the contract of
\citep[Ch.~6.2]{iiyama2026thesis} and \citep{iiyama2026scud}; the corresponding
thesis equation numbers are cited inline so that the two documents can be read
side by side.

Every filter in this chapter is a \emph{discrete-time, continuous-dynamics}
estimator: the state is propagated by numerical integration of the dynamics
(\cref{chap:dynamics,chap:integration}), the state-transition matrix and
measurement Jacobian are produced by automatic differentiation
(\cref{chap:autodiff}), and the covariance algebra specified below is the only
part that distinguishes one filter from another. The algorithms therefore share
a single interface, a single outlier-rejection policy, and a single logging
contract; they differ only in \emph{how} the second moment is represented and
propagated, and hence in their numerical conditioning and cost.

\section{The Recursive Estimation Contract}
\label{sec:flt-contract}

\subsection{State and covariance}

Every recursive filter derives from the abstract base class \cpp{Filter}
(\file{cpp/lupnt/numerics/filters/filter.h}) and maintains a state estimate and
a covariance
\begin{equation}
  \label{eq:flt-state}
  \hat{\vec{x}} \in \mathbb{R}^{n},
  \qquad
  \mat{P} \in \mathbb{R}^{n\times n},
  \qquad
  \mat{P} = \mat{P}\transpose \succeq 0 ,
\end{equation}
where $n$ is the state dimension. The base class keeps three snapshots of the
pair: the working copy $(\hat{\vec{x}}, \mat{P})$ in \cpp{x\_}/\cpp{P\_}, the
prior (post-propagation, pre-update) copy
$(\hat{\vec{x}}^{-}, \mat{P}^{-})$ in \cpp{x\_prior\_}/\cpp{P\_prior\_}, and
the posterior (post-update) copy $(\hat{\vec{x}}^{+}, \mat{P}^{+})$ in
\cpp{x\_post\_}/\cpp{P\_post\_}. A fourth matrix, \cpp{P\_bar\_}, holds the
\emph{propagated-only} covariance
\begin{equation}
  \label{eq:flt-pbar}
  \bar{\mat{P}} = \mat{F}\,\mat{P}\,\mat{F}\transpose ,
\end{equation}
that is, the covariance after the dynamics but before process noise is added.
It is exposed through \cpp{GetCovarianceBar} and is consumed only by the
adaptive process-noise estimator of \cref{sec:flt-adaptive}.

The state is carried as a \cpp{State} object, which attaches a reference frame
and a unit convention to the raw vector; the covariance is a plain
\cpp{MatXd}. It is the caller's responsibility to keep the two consistent: the
covariance is always expressed in the same units and frame as the state it
accompanies, and no filter performs an implicit unit conversion. For a coupled
orbit--clock state the usual layout is
\begin{equation}
  \label{eq:flt-layout}
  \vec{x} =
  \begin{bmatrix}
    \vec{r}\transpose & \vec{v}\transpose & b & \dot{b} & (\ddot{b})
  \end{bmatrix}\transpose ,
\end{equation}
with $\vec{r}, \vec{v} \in \mathbb{R}^3$ the position and velocity in the
configured frame and $(b, \dot{b}, \ddot{b})$ the clock bias, drift, and
drift rate in the configured clock-bias unit
(\cref{sec:flt-clock-noise}).

\subsection{The three callbacks}

A filter is not a model. It is driven entirely by three callbacks registered
before the estimation loop starts; all model knowledge enters through them. The
first is the dynamics callback,
\begin{equation}
  \label{eq:flt-cb-dyn}
  \texttt{f\_dyn\_} : (\vec{x}, t_0, t_f, \vec{u}) \;\longmapsto\;
  \bigl(\vec{x}(t_f),\ \mat{F}\bigr),
  \qquad
  \mat{F} = \frac{\partial \vec{x}(t_f)}{\partial \vec{x}(t_0)} ,
\end{equation}
which propagates the state over $[t_0, t_f]$ and returns the state-transition
matrix (STM) $\mat{F}$ of that propagation. The second is the process-noise
callback,
\begin{equation}
  \label{eq:flt-cb-proc}
  \texttt{f\_proc\_} : (\vec{x}, t_0, t_f) \;\longmapsto\; \mat{Q},
  \qquad
  \mat{Q} = \mat{Q}\transpose \succeq 0 ,
\end{equation}
which returns the process-noise covariance accumulated over the same interval.
The third is the measurement callback,
\begin{equation}
  \label{eq:flt-cb-meas}
  \texttt{f\_meas\_} : \vec{x} \;\longmapsto\;
  \bigl(\hat{\vec{z}},\ \mat{H},\ \mat{R}\bigr),
  \qquad
  \mat{H} = \frac{\partial \vec{z}}{\partial \vec{x}} ,
\end{equation}
which predicts the measurement at the current estimate and returns its Jacobian
and noise covariance. The callback signatures are the typedefs
\cpp{FilterDynamicsFunction}, \cpp{ProcessNoiseFunction}, and
\cpp{FilterMeasurementFunction} declared in \file{filter.h}.

Plain models (a \cpp{DynamicsFunction} or a \cpp{MeasurementFunction} without a
Jacobian) are lifted into the callback forms of \cref{eq:flt-cb-dyn} and
\cref{eq:flt-cb-meas} by \cpp{GetFilterDynamicsFunction} and
\cpp{GetFilterMeasurementFunction}
(\file{cpp/lupnt/numerics/filters/filter_utils.h}), which obtain $\mat{F}$ and
$\mat{H}$ by forward-mode automatic differentiation of the model
(\cref{chap:autodiff}). No filter ever forms a finite-difference Jacobian.

\begin{lupntnote}
The process noise is evaluated \emph{at the prior state}, before the dynamics
are applied: $\mat{Q} = \texttt{f\_proc\_}(\hat{\vec{x}}, t_0, t_f)$. All
filters follow this convention so that a state-dependent $\mat{Q}$ (for example
one keyed on the current altitude) is reproducible across the family.
\end{lupntnote}

\subsection{The predict--update cycle}

The estimation loop alternates two virtual calls,
\begin{equation}
  \label{eq:flt-cycle}
  \cdots \;\longrightarrow\;
  \underbrace{\texttt{Predict}(t_k)}_{(\hat{\vec{x}}^{+}_{k-1},\mat{P}^{+}_{k-1})
    \,\mapsto\,(\hat{\vec{x}}^{-}_{k},\mat{P}^{-}_{k})}
  \;\longrightarrow\;
  \underbrace{\texttt{Update}(\vec{z}_k)}_{(\hat{\vec{x}}^{-}_{k},\mat{P}^{-}_{k})
    \,\mapsto\,(\hat{\vec{x}}^{+}_{k},\mat{P}^{+}_{k})}
  \;\longrightarrow\; \cdots
\end{equation}
Both are no-ops for an empty measurement vector: if $\vec{z}_k$ has zero
length, \cpp{Update} returns immediately and the prior becomes the posterior.
This is the normal path for an epoch with no visible transmitters.

\subsection{Process-noise mapping}

Correlated process noise on a subset of the state is supplied through a
mapping matrix $\mat{G} \in \mathbb{R}^{n\times m}$ registered by
\cpp{SetProcessNoiseMappingMatrix}. When mapping is enabled the noise is
transformed before it is used,
\begin{equation}
  \label{eq:flt-gmap}
  \mat{Q} \;\longleftarrow\; \mat{G}\,\mat{Q}\,\mat{G}\transpose ,
\end{equation}
so that the callback of \cref{eq:flt-cb-proc} may return a compact
$m\times m$ noise covariance in the natural noise coordinates (for example a
three-axis acceleration PSD) and the filter expands it to the full state. For
the dense filters \cref{eq:flt-gmap} is applied literally; for the factorized
filters of \cref{sec:flt-udu,sec:flt-srif} the mapping is kept
\emph{separate} from the state factor and never multiplied out, which is
precisely what preserves the small variances (see \cref{sec:flt-udu}).

\subsection{Outlier rejection and fault detection}
\label{sec:flt-outliers}

Before any gain is applied, every measurement component is screened by a
normalized-residual test. With the innovation covariance
$\mat{S} = \mat{R} + \mat{H}\mat{P}^{-}\mat{H}\transpose$ and the residual
$\delta \vec{z} = \vec{z} - \hat{\vec{z}}$, component $i$ is rejected when
\begin{equation}
  \label{eq:flt-outlier}
  \frac{\lvert \delta z_i \rvert}{\sqrt{S_{ii}}} \;>\; \tau ,
\end{equation}
where $\tau$ is \cpp{outlier\_threshold\_}, default $\tau = 3$ (a
$3\sigma$ gate) and configurable from YAML through the
\code{outlier\_threshold} key or at run time through
\cpp{SetOutlierThreshold}. The flagged rows are dropped from $\delta \vec{z}$,
$\mat{H}$, and both the rows and columns of $\mat{R}$, and the innovation
covariance is recomputed on the shrunken system,
\begin{equation}
  \label{eq:flt-outlier-recompute}
  \mat{S} \;\longleftarrow\; \mat{H}\,\mat{P}^{-}\,\mat{H}\transpose + \mat{R} .
\end{equation}
Note that \cref{eq:flt-outlier-recompute} uses the stored \emph{prior}
covariance \cpp{P\_prior\_}, not the working copy, so the recomputation is
exact rather than a rank-corrected approximation. If every component is
rejected, the update returns without touching the estimate.

\begin{lstlisting}[language=cpp, caption={Normalized-residual outlier gate.},
                   label={lst:flt-outliers}]
// cpp/lupnt/numerics/filters/ekf.cc :: EKF::RemoveOutliers
VecXd ratio = dz_.array().abs() / S_.diagonal().array().sqrt();
VecXb is_outlier = ratio.array() > outlier_threshold_;

int N_meas_new = N_meas - is_outlier.count();
if (N_meas_new == N_meas) return;
// ... copy the surviving rows of dz_, H_ and the surviving rows/cols of R_ ...
H_ = H_new;
R_ = R_new;
dz_ = dz_new;
S_ = H_ * P_prior_ * H_.transpose() + R_;   // exact, on the shrunken system
\end{lstlisting}

The default gate can be replaced wholesale by registering a
\cpp{FaultDetectionFunction} through \cpp{SetFaultDetectionFunction}. The
callback has signature \cpp{void(KalmanFilter*)} and is handed the filter
itself, so it may inspect and rewrite $\delta \vec{z}$, $\mat{H}$, $\mat{R}$,
and $\mat{S}$ by any rule --- a chi-square test on the full innovation vector,
a RAIM-style subset-exclusion search, or a clock-fault monitor built on
inter-satellite ranges \citep{iiyama2026edm}. Once a custom function is
registered, \cref{eq:flt-outlier} is not evaluated at all. The same screening
code is shared by \cpp{EKF}, \cpp{SRIF}, and \cpp{UDUEKF}; the UKF does not
screen, because it never forms a row-wise Jacobian.

\begin{lupntwarning}
The gate of \cref{eq:flt-outlier} is a per-component test against the
\emph{filter's own} innovation covariance. A slowly diverging filter inflates
$S_{ii}$, passes increasingly wrong measurements, and can lock onto a false
solution. For safety-relevant use, register a fault-detection function with an
independent consistency statistic rather than relying on the default gate.
\end{lupntwarning}

\section{Extended Kalman Filter}
\label{sec:flt-ekf}

The \cpp{EKF} is the reference implementation of the family; every other dense
filter inherits its screening, logging, and smoothing machinery.

\subsection{Time update}

\cpp{EKF::Predict} evaluates the process noise at the prior state, propagates
the state and STM through \cref{eq:flt-cb-dyn}, and forms the propagated
covariance:
\begin{align}
  \label{eq:flt-ekf-q}
  \mat{Q} &= \texttt{f\_proc\_}(\hat{\vec{x}}, t_0, t_f),
  \qquad
  \mat{Q} \leftarrow \mat{G}\mat{Q}\mat{G}\transpose \ \text{(if mapped)}, \\
  \label{eq:flt-ekf-x}
  \hat{\vec{x}}^{-} &= \texttt{f\_dyn\_}(\hat{\vec{x}}, t_0, t_f, \vec{u};\, \mat{F}), \\
  \label{eq:flt-ekf-p}
  \bar{\mat{P}} &= \mat{F}\,\mat{P}\,\mat{F}\transpose,
  \qquad
  \mat{P}^{-} = \bar{\mat{P}} + \mat{Q} .
\end{align}
The linearization of \cref{eq:flt-ekf-p} is the defining approximation of the
extended filter: the covariance is transported by the Jacobian of the true
nonlinear flow evaluated at the current estimate, so the propagated second
moment is correct to first order in the state error.

\begin{lstlisting}[language=cpp, caption={EKF time update.},
                   label={lst:flt-ekf-predict}]
// cpp/lupnt/numerics/filters/ekf.cc :: EKF::Predict
Q_ = f_proc_(x_, t_, t);
if (use_process_noise_mapping_) Q_ = G_ * Q_ * G_.transpose();

x_ = f_dyn_(x_, t_, t, u, &F_);

P_bar_ = F_ * P_ * F_.transpose();   // retained for adaptive process noise
P_ = P_bar_ + Q_;

t_ = t;
x_prior_ = x_;
P_prior_ = P_;
\end{lstlisting}

\subsection{Measurement update}

\cpp{EKF::Update} linearizes the measurement about the prior, screens the
residual (\cref{sec:flt-outliers}), and applies a Joseph-form correction:
\begin{align}
  \label{eq:flt-ekf-innov}
  \hat{\vec{z}} &= \texttt{f\_meas\_}(\hat{\vec{x}}^{-};\, \mat{H}, \mat{R}),
  \qquad
  \mat{S} = \mat{R} + \mat{H}\,\mat{P}^{-}\,\mat{H}\transpose,
  \qquad
  \delta \vec{z} = \vec{z} - \hat{\vec{z}}, \\
  \label{eq:flt-ekf-gain}
  \mat{K} &= \mat{P}^{-}\mat{H}\transpose \mat{S}^{-1},
  \qquad
  \delta \vec{x} = \mat{K}\,\delta \vec{z},
  \qquad
  \hat{\vec{x}}^{+} = \hat{\vec{x}}^{-} + \delta \vec{x}, \\
  \label{eq:flt-ekf-joseph}
  \mat{P}^{+} &= (\mat{I} - \mat{K}\mat{H})\,\mat{P}^{-}\,
                (\mat{I} - \mat{K}\mat{H})\transpose
              + \mat{K}\,\mat{R}\,\mat{K}\transpose .
\end{align}
The covariance of the state correction itself,
\begin{equation}
  \label{eq:flt-ekf-sigmadx}
  \mat{\Sigma}_{\delta x} = \mat{K}\,\mat{S}\,\mat{K}\transpose ,
\end{equation}
is recorded for the adaptive process-noise estimator of
\cref{sec:flt-adaptive}.

\Cref{eq:flt-ekf-joseph} is the Joseph (symmetric, or ``stabilized'') form
rather than the shorter $\mat{P}^{+} = (\mat{I} - \mat{K}\mat{H})\mat{P}^{-}$.
Two properties motivate the extra work. First, it is manifestly symmetric and
is a sum of two positive-semidefinite congruences, so round-off cannot
introduce an asymmetric or indefinite result to first order. Second --- and
this is what the Schmidt filter below depends on --- it is valid for
\emph{any} gain $\mat{K}$, not only the optimal one: expanding the covariance
of $\hat{\vec{x}}^{+} = \hat{\vec{x}}^{-} + \mat{K}(\vec{z} - \mat{H}
\hat{\vec{x}}^{-})$ for arbitrary $\mat{K}$ reproduces
\cref{eq:flt-ekf-joseph} exactly, whereas the short form is only correct when
$\mat{K}$ satisfies the normal equations. Its cost is $O(n^3)$ per update
against $O(n^2 m)$ for the short form, which is negligible at the state sizes
relevant here.

\begin{lstlisting}[language=cpp, caption={EKF measurement update, Joseph form.},
                   label={lst:flt-ekf-update}]
// cpp/lupnt/numerics/filters/ekf.cc :: EKF::Update
z_prior_ = f_meas_(x_, &H_, &R_);
S_ = R_ + H_ * P_ * H_.transpose();
dz_ = z_true_ - z_prior_;

if (use_custom_fault_detection_ && f_fault_det_) f_fault_det_(this);
else                                             RemoveOutliers();
if (dz_.size() == 0) return;                 // all measurements rejected

K_ = P_ * H_.transpose() * S_.inverse();     // Kalman gain
if (n_consider_ > 0) K_.bottomRows(n_consider_).setZero();

dx_ = K_ * dz_;
Sigma_dx_ = K_ * S_ * K_.transpose();
x_ = x_ + dx_;

MatXd I = MatXd::Identity(N_state, N_state);
MatXd G = I - K_ * H_;
P_ = G * P_ * G.transpose() + K_ * R_ * K_.transpose();   // Joseph form
\end{lstlisting}

\subsection{Consider (Schmidt) states}
\label{sec:flt-schmidt}

Calling \cpp{SetConsiderStateCount}$(n_c)$ --- also exposed as the thin
subclass \cpp{SchmidtEKF} in \file{schmidt_ekf.h} --- marks the trailing $n_c$
components of the state as \emph{consider parameters}. Partition
\begin{equation}
  \label{eq:flt-schmidt-part}
  \vec{x} = \begin{bmatrix} \vec{x}_s \\ \vec{x}_c \end{bmatrix},
  \qquad
  \mat{P} = \begin{bmatrix}
    \mat{P}_{ss} & \mat{P}_{sc} \\
    \mat{P}_{sc}\transpose & \mat{P}_{cc}
  \end{bmatrix},
  \qquad
  \vec{x}_c \in \mathbb{R}^{n_c} .
\end{equation}
Consider parameters enter $\mat{F}$, $\mat{H}$, $\mat{K}$, and the covariance
bookkeeping in the usual way, but their mean is never corrected: the trailing
block of the gain is zeroed,
\begin{equation}
  \label{eq:flt-schmidt-gain}
  \mat{K} =
  \begin{bmatrix} \mat{K}_s \\ \mat{K}_c \end{bmatrix}
  \;\longrightarrow\;
  \tilde{\mat{K}} =
  \begin{bmatrix} \mat{K}_s \\ \mat{0} \end{bmatrix} ,
\end{equation}
before $\delta \vec{x} = \tilde{\mat{K}}\,\delta \vec{z}$ is formed. Because
\cref{eq:flt-ekf-joseph} holds for the suboptimal gain $\tilde{\mat{K}}$, the
consider block's variance $\mat{P}_{cc}$ and, crucially, its cross-covariance
$\mat{P}_{sc}$ with the solved-for states still update consistently. The effect
is that the uncertainty of an unestimated bias (an ephemeris error, an
unmodelled empirical acceleration, a station coordinate) is correctly
propagated into the estimated states without the filter attempting --- and, for
a weakly observable parameter, failing --- to solve for it.

\begin{lstlisting}[language=cpp, caption={Schmidt/consider gain truncation.},
                   label={lst:flt-schmidt}]
// cpp/lupnt/numerics/filters/ekf.cc :: EKF::Update (n_consider_ branch)
if (n_consider_ > 0) {
  // Schmidt-Kalman: never correct the trailing consider states. The Joseph-form
  // covariance update remains valid for this non-optimal gain, so the consider
  // states' (co)variance still updates consistently even though their mean does
  // not move.
  K_.bottomRows(n_consider_).setZero();
}
dx_ = K_ * dz_;
\end{lstlisting}

\section{Unscented Kalman Filter}
\label{sec:flt-ukf}

\cpp{UKF} replaces the analytic linearization of \cref{eq:flt-ekf-p} and
\cref{eq:flt-ekf-innov} with a deterministic sampling scheme, the unscented
transform of \citet{julier2004ukf}. Rather than propagating a Jacobian, it
propagates a small, carefully chosen set of points whose sample mean and
covariance match those of the prior, and reconstructs the moments of the
transformed distribution from the transformed points. The resulting moments are
accurate to second order in the Taylor expansion of the nonlinearity for any
distribution, and to third order for a Gaussian --- against first order for the
EKF --- at the price of $2n+1$ dynamics evaluations per step.

\subsection{Scaling and weights}

For an $n$-dimensional state the filter uses $2n+1$ sigma points with the
composite scaling parameter
\begin{equation}
  \label{eq:flt-ukf-lambda}
  \lambda = \alpha^2 (n + \kappa) - n ,
\end{equation}
and the mean and covariance weights
\begin{equation}
  \label{eq:flt-ukf-weights}
  w^m_0 = \frac{\lambda}{n+\lambda},
  \qquad
  w^c_0 = \frac{\lambda}{n+\lambda} + \bigl(1 - \alpha^2 + \beta\bigr),
  \qquad
  w^m_i = w^c_i = \frac{1}{2(n+\lambda)},
  \quad i = 1,\dots,2n .
\end{equation}
The three tuning constants have distinct roles, summarized in
\cref{tab:flt-ukf-params}: $\alpha$ sets the spread of the points about the
mean and is normally small so that the sampling stays in the region where the
local linearization is meaningful; $\kappa$ is a secondary scaling that shifts
weight between the centre point and the symmetric pairs; $\beta$ injects prior
knowledge of the fourth moment, with $\beta = 2$ optimal for a Gaussian prior.
Note that $\sum_i w^m_i = \sum_i w^c_i = 1$ by construction, and that $w^m_0$
and $w^c_0$ are negative whenever $\lambda < 0$, which is the normal regime for
$\alpha \ll 1$. A negative centre weight is legitimate --- the reconstructed
covariance remains positive semidefinite for the standard parameter ranges ---
but it is the mechanism by which an ill-conditioned prior can produce a
non-positive reconstructed covariance and abort the next Cholesky factorization.

\begin{table}[H]
  \centering
  \caption{Unscented-transform tuning parameters and their defaults in
           \file{cpp/lupnt/numerics/filters/ukf.h}.}
  \label{tab:flt-ukf-params}
  \begin{tabularx}{\textwidth}{llL}
    \toprule
    Symbol & Member (default) & Role \\
    \midrule
    $\alpha$ & \cpp{alpha\_} ($10^{-3}$) &
      Spread of the sigma points about the mean; small values keep the sampling
      local. Set with \cpp{SetAlpha}. \\
    $\beta$ & \cpp{beta\_} ($2$) &
      Prior knowledge of the distribution's kurtosis; $\beta = 2$ is optimal
      for a Gaussian. Enters $w^c_0$ only. Set with \cpp{SetBeta}. \\
    $\kappa$ & \cpp{kappa\_} ($0$) &
      Secondary scaling; $\kappa = 3 - n$ matches the fourth moment of a
      Gaussian but yields a negative $\lambda$ for $n > 3$. Set with
      \cpp{SetKappa}. \\
    $\lambda$ & \cpp{lambda\_} (derived) &
      Composite scaling of \cref{eq:flt-ukf-lambda}. \\
    \bottomrule
  \end{tabularx}
\end{table}

\begin{lstlisting}[language=cpp, caption={Unscented-transform weights.},
                   label={lst:flt-ukf-weights}]
// cpp/lupnt/numerics/filters/ukf.cc :: UKF::InitializeUkfParams
lambda_  = alpha_ * alpha_ * (n_x_ + kappa_) - n_x_;
n_sigma_ = 2 * n_x_ + 1;

w_m_(0) = lambda_ / (n_x_ + lambda_);
w_c_(0) = lambda_ / (n_x_ + lambda_) + (1.0 - alpha_ * alpha_ + beta_);

for (int i = 1; i < n_sigma_; i++) {
  w_m_(i) = 1.0 / (2.0 * (n_x_ + lambda_));
  w_c_(i) = 1.0 / (2.0 * (n_x_ + lambda_));
}
\end{lstlisting}

\subsection{Sigma points}

The points about $(\hat{\vec{x}}, \mat{P})$ are
\begin{equation}
  \label{eq:flt-ukf-sigma}
  \mathcal{X}_0 = \hat{\vec{x}},
  \qquad
  \mathcal{X}_i = \hat{\vec{x}} + \bigl[\mat{L}\bigr]_i,
  \qquad
  \mathcal{X}_{i+n} = \hat{\vec{x}} - \bigl[\mat{L}\bigr]_i,
  \qquad i = 1,\dots,n ,
\end{equation}
where $[\mat{L}]_i$ is the $i$-th column of the lower-triangular Cholesky
factor of the scaled covariance,
\begin{equation}
  \label{eq:flt-ukf-chol}
  \mat{L}\,\mat{L}\transpose = (n + \lambda)\,\mat{P} .
\end{equation}
The set $\{\mathcal{X}_i\}$ with the weights of \cref{eq:flt-ukf-weights}
reproduces the first two moments of the prior exactly.

\begin{lstlisting}[language=cpp, caption={Sigma-point generation.},
                   label={lst:flt-ukf-sigma}]
// cpp/lupnt/numerics/filters/ukf.cc :: UKF::ComputeSigmaPoints
MatXd chol_cov = ((n_x_ + lambda_) * cov).llt().matrixL();

MatX sigma_points(n_x_, n_sigma_);
sigma_points.col(0) = state;
for (int i = 0; i < n_x_; i++) {
  sigma_points.col(i + 1)        = state + chol_cov.col(i);
  sigma_points.col(i + 1 + n_x_) = state - chol_cov.col(i);
}
\end{lstlisting}

\subsection{Time and measurement updates}

\cpp{UKF::Predict} propagates each point through \texttt{f\_dyn\_} with a null
Jacobian pointer --- no STM is computed --- and recombines:
\begin{align}
  \label{eq:flt-ukf-prop}
  \mathcal{X}_i^{-} &= \texttt{f\_dyn\_}(\mathcal{X}_i, t_0, t_f, \vec{u};\, \text{nullptr}), \\
  \label{eq:flt-ukf-predict}
  \hat{\vec{x}}^{-} &= \sum_{i=0}^{2n} w^m_i\,\mathcal{X}_i^{-},
  \qquad
  \mat{P}^{-} = \sum_{i=0}^{2n} w^c_i
      \bigl(\mathcal{X}_i^{-} - \hat{\vec{x}}^{-}\bigr)
      \bigl(\mathcal{X}_i^{-} - \hat{\vec{x}}^{-}\bigr)\transpose + \mat{Q} .
\end{align}
\cpp{UKF::Update} regenerates the points at the prior, maps them through
\texttt{f\_meas\_} to measurement points $\mathcal{Z}_i$, and forms the
predicted measurement, its covariance, and the state--measurement
cross-covariance:
\begin{align}
  \label{eq:flt-ukf-meas}
  \hat{\vec{z}} &= \sum_{i=0}^{2n} w^m_i\,\mathcal{Z}_i,
  \qquad
  \mat{S} = \sum_{i=0}^{2n} w^c_i
      \bigl(\mathcal{Z}_i - \hat{\vec{z}}\bigr)
      \bigl(\mathcal{Z}_i - \hat{\vec{z}}\bigr)\transpose + \mat{R}, \\
  \label{eq:flt-ukf-cross}
  \mat{P}_{xz} &= \sum_{i=0}^{2n} w^c_i
      \bigl(\mathcal{X}_i - \hat{\vec{x}}^{-}\bigr)
      \bigl(\mathcal{Z}_i - \hat{\vec{z}}\bigr)\transpose,
  \qquad
  \mat{K} = \mat{P}_{xz}\,\mat{S}^{-1}, \\
  \label{eq:flt-ukf-update}
  \hat{\vec{x}}^{+} &= \hat{\vec{x}}^{-} + \mat{K}\bigl(\vec{z} - \hat{\vec{z}}\bigr),
  \qquad
  \mat{P}^{+} = \mat{P}^{-} - \mat{K}\,\mat{S}\,\mat{K}\transpose .
\end{align}
\Cref{eq:flt-ukf-cross} is the unscented replacement for
$\mat{P}^{-}\mat{H}\transpose$: the gain is defined directly by the sampled
cross-covariance, so no measurement Jacobian is ever formed. The measurement
noise $\mat{R}$ used in \cref{eq:flt-ukf-meas} is the one returned at the first
sigma point; the implementation evaluates $\mat{R}$ at every point but keeps
only $\mat{R}_0$, which is exact whenever $\mat{R}$ is state-independent and a
first-order approximation otherwise.

\begin{lupntwarning}
\Cref{eq:flt-ukf-update} uses the short covariance form, not the Joseph form.
The UKF is therefore the least numerically robust member of the family: a long
run with tightly constrained states can drive $\mat{P}$ indefinite and abort
the Cholesky factorization of \cref{eq:flt-ukf-chol}. The UKF also does not run
the outlier gate of \cref{eq:flt-outlier} and does not maintain the smoother
logs. Use it when the nonlinearity, not the conditioning, is the limiting
factor.
\end{lupntwarning}

\section{UD-Factorized Filter (Bierman--Thornton)}
\label{sec:flt-udu}

\cpp{UDUEKF} never stores a covariance matrix as such. It maintains the
factorization
\begin{equation}
  \label{eq:flt-udu-fact}
  \mat{P} = \mat{U}\,\diag(\vec{D})\,\mat{U}\transpose ,
\end{equation}
with $\mat{U}$ unit upper-triangular ($U_{ii} = 1$, $U_{ij} = 0$ for $i > j$)
and $\vec{D}$ a non-negative diagonal. This is the UD factorization of
\citet{bierman1977} and \citet{thornton1980}, and it mirrors the thesis
factorization
$\hat{\mat{P}}_{k|m} = \mat{U}_{k|m}\mat{D}_{k|m}\mat{U}_{k|m}\transpose$
\citep[Ch.~6.2.1, Eq.~6.17, Algorithm~1]{iiyama2026thesis}.

Three properties follow immediately. Symmetry and positive semidefiniteness are
structural, not numerical: any $(\mat{U}, \vec{D})$ with $\vec{D} \geq 0$
reconstructs a valid covariance, so no amount of round-off can produce an
indefinite $\mat{P}$. The effective condition number of the stored quantity is
the square root of that of $\mat{P}$, so the factored filter carries roughly
twice the significant digits of the dense filter for the same word length.
Finally --- and this is the decisive property for coupled orbit--clock
estimation --- variances many orders of magnitude apart survive. A lunar
orbit--clock state routinely spans from $\sim\!\SI{1e2}{\meter\squared}$ on
position down to $\sim\!10^{-26}$ on clock drift; forming
$\mat{P} = \mat{U}\diag(\vec{D})\mat{U}\transpose$, adding $\mat{Q}$, and
re-factoring rounds the smallest diagonal entries to zero within a few steps
and collapses the drift estimate. The algorithms below are constructed so that
$\mat{P}$ is never re-formed inside the recursion.

The initial factorization is obtained once, from the user-supplied covariance,
by a backward Cholesky-like sweep in \cpp{UDUDecomposition}
(\file{cpp/lupnt/numerics/filters/udu_utils.cc}):
\begin{equation}
  \label{eq:flt-udu-decomp}
  D_j = P_{jj} - \sum_{k>j} U_{jk}^2 D_k ,
  \qquad
  U_{ij} = \frac{1}{D_j}\Bigl(P_{ij} - \sum_{k>j} U_{ik} D_k U_{jk}\Bigr),
  \quad i < j ,
\end{equation}
evaluated for $j = n, n-1, \dots, 1$; a $D_j$ below $10^{-16}$ in magnitude is
clamped to zero and its column of $\mat{U}$ left at the identity.

\subsection{Time update: the Thornton MWGS}

The predict step factors the propagated covariance directly. Define the
$n \times (n+m)$ coefficient array and the $(n+m)\times(n+m)$ weight matrix
\begin{equation}
  \label{eq:flt-udu-mwgs-y}
  \mat{Y} = \begin{bmatrix} \mat{F}\mat{U} & \mat{G} \end{bmatrix},
  \qquad
  \tilde{\mat{D}} = \operatorname{blkdiag}\bigl(\diag(\vec{D}),\ \mat{Q}\bigr) ,
\end{equation}
with $\mat{G} \in \mathbb{R}^{n\times m}$ the process-noise mapping of
\cref{eq:flt-gmap} (identity when no mapping is registered) and $\mat{Q}$
diagonal. A modified weighted Gram--Schmidt (MWGS) orthogonalization of the
rows of $\mat{Y}$ in the $\tilde{\mat{D}}$ inner product returns
$(\mat{U}^{-}, \vec{D}^{-})$ such that
\begin{equation}
  \label{eq:flt-udu-mwgs}
  \mat{U}^{-}\diag(\vec{D}^{-})\,(\mat{U}^{-})\transpose
  = \mat{Y}\,\tilde{\mat{D}}\,\mat{Y}\transpose
  = \mat{F}\,\mat{P}\,\mat{F}\transpose + \mat{G}\,\mat{Q}\,\mat{G}\transpose ,
\end{equation}
where the right-hand product is \emph{never formed}
\citep[Ch.~6.2.2, Eqs.~6.20--6.23]{iiyama2026thesis}. Writing
$\vec{b}_j$ for the $j$-th row of $\mat{Y}$, the recursion sweeps
$j = n, n-1, \dots, 1$ (Thornton ordering) and computes
\begin{equation}
  \label{eq:flt-udu-mwgs-rec}
  D^{-}_j = \vec{b}_j\transpose \tilde{\mat{D}}\,\vec{b}_j ,
  \qquad
  U^{-}_{ij} = \frac{\vec{b}_i\transpose \tilde{\mat{D}}\,\vec{b}_j}{D^{-}_j},
  \qquad
  \vec{b}_i \leftarrow \vec{b}_i - U^{-}_{ij}\,\vec{b}_j ,
  \quad i < j .
\end{equation}
Only an exactly non-positive $D^{-}_j$ is clamped to zero; a tiny but positive
variance is deliberately kept, because that is what a clock-drift term looks
like.

\begin{lstlisting}[language=cpp, caption={Thornton MWGS time update of the UD factors.},
                   label={lst:flt-udu-predict}]
// cpp/lupnt/numerics/filters/udu_filter.cc :: UDUEKF::Predict
Q_ = f_proc_(x_, t_, t);
x_ = f_dyn_(x_, t_, t, u, &F_);

if (!use_process_noise_mapping_) G_ = MatXd::Identity(n, n);
LUPNT_CHECK(Q_.isDiagonal(), "UDU process noise must be diagonal ...", "UDU");
const int m = static_cast<int>(G_.cols());   // reduced process-noise dimension

P_bar_ = F_ * UDUReconstruct(U_, D_diag_) * F_.transpose();  // F P F^T (adaptive only)

MatXd Y = MatXd::Zero(n, n + m);
Y.leftCols(n)  = F_ * U_;
Y.rightCols(m) = G_;
MatXd D_tilde = MatXd::Zero(n + m, n + m);
D_tilde.topLeftCorner(n, n)     = D_diag_.asDiagonal();
D_tilde.bottomRightCorner(m, m) = Q_;
std::tie(D_diag_, U_) = ModifiedGramSchmidt(D_tilde, Y);
\end{lstlisting}

\begin{lstlisting}[language=cpp, caption={The MWGS kernel.},
                   label={lst:flt-udu-mwgs}]
// cpp/lupnt/numerics/filters/udu_utils.cc :: ModifiedGramSchmidt
std::vector<VecXd> b(n);
for (int i = 0; i < n; ++i) b[i] = Y.row(i).transpose();

for (int k = 0; k < n; ++k) {
  const int j = (n - 1) - k;           // Thornton ordering: last row first
  VecXd fj = D_tilde * b[j];
  D_m(j) = b[j].dot(fj);
  if (D_m(j) <= 0.0) { D_m(j) = 0.0; continue; }
  fj /= D_m(j);
  for (int i = 0; i < j; ++i) {
    U_m(i, j) = b[i].dot(fj);
    b[i] -= U_m(i, j) * b[j];
  }
}
\end{lstlisting}

\subsection{Measurement update: the Carlson rank-one recursion}

Measurements are assimilated \emph{sequentially}, one scalar component at a
time, which requires a diagonal $\mat{R}$. For a scalar row
$\vec{h}\transpose$ with variance $r$, define
\begin{equation}
  \label{eq:flt-udu-carlson-f}
  \vec{f} = \mat{U}\transpose \vec{h},
  \qquad
  \vec{v} = \diag(\vec{D})\,\vec{f},
  \qquad\text{so that}\qquad
  \vec{f}\transpose \vec{v} = \vec{h}\transpose \mat{P}\,\vec{h} .
\end{equation}
The recursion accumulates the innovation variance,
\begin{equation}
  \label{eq:flt-udu-alpha}
  \alpha_0 = r,
  \qquad
  \alpha_k = \alpha_{k-1} + v_k f_k,
  \qquad k = 1, \dots, n ,
\end{equation}
so that $\alpha_n = r + \vec{h}\transpose\mat{P}\vec{h} = S$; because
$\alpha_0 = r > 0$ and each increment $v_k f_k = D_k f_k^2 \geq 0$, the
sequence is positive and increasing, and the recursion cannot divide by zero.
The diagonal is downdated by the ratio of consecutive $\alpha$'s,
\begin{equation}
  \label{eq:flt-udu-d}
  D^{+}_k = \frac{\alpha_{k-1}}{\alpha_k}\,D_k \;\leq\; D_k ,
\end{equation}
which is manifestly a contraction, and the columns of $\mat{U}$ receive a
rank-one correction driven by the running gain vector $\vec{\beta}$,
\begin{equation}
  \label{eq:flt-udu-u}
  \beta_k \leftarrow v_k,
  \qquad
  p_k = -\frac{f_k}{\alpha_{k-1}},
  \qquad
  \begin{aligned}
    U^{+}_{jk} &= U_{jk} + \beta_j\,p_k, \\
    \beta_j    &\leftarrow \beta_j + U_{jk}\,v_k,
  \end{aligned}
  \qquad j < k .
\end{equation}
The Kalman gain for the scalar measurement falls out of the same accumulator,
\begin{equation}
  \label{eq:flt-udu-gain}
  \vec{k} = \frac{\vec{\beta}}{\alpha_n} ,
\end{equation}
which is the rank-one UD downdate of
\citep[Ch.~6.2.3, Eqs.~6.24--6.27, Algorithm~2]{iiyama2026thesis}. Because the
components are processed one at a time against a state that has already
absorbed the previous ones, the residual of the $m$-th component must be
corrected for the accumulated correction,
\begin{equation}
  \label{eq:flt-udu-residual}
  \delta z_m' = \delta z_m - \vec{h}_m\transpose \delta \vec{x},
  \qquad
  \delta \vec{x} \leftarrow \delta \vec{x} + \vec{k}\,\delta z_m' ,
\end{equation}
which reproduces the batch update exactly for a linear model. The dense
$\mat{P}$ is reconstructed once, at the end of the update, purely for
reporting.

\begin{lstlisting}[language=cpp, caption={Carlson rank-one sequential UD update.},
                   label={lst:flt-udu-carlson}]
// cpp/lupnt/numerics/filters/udu_filter.cc :: UDUEKF::CarlsonUpdate
for (int meas = 0; meas < n_meas; ++meas) {
  VecXd h = H_.row(meas).transpose();
  double r = R_(meas, meas);

  VecXd f = U_.transpose() * h;        // f = U^T h
  VecXd v = D_diag_.asDiagonal() * f;  // v = D f, so f.dot(v) = h^T P h

  alpha(0) = r;
  for (int k = 0; k < n; ++k) {
    alpha(k + 1) = alpha(k) + v(k) * f(k);  // grows from r > 0, stays positive
    D_plus(k) = alpha(k) / alpha(k + 1) * D_diag_(k);
    beta(k) = v(k);
    if (k == 0) continue;
    double p_k = -f(k) / alpha(k);
    for (int j = 0; j < k; ++j) {
      U_plus(j, k) = U_(j, k) + beta(j) * p_k;
      beta(j)      = beta(j) + U_(j, k) * v(k);
    }
  }

  VecXd k_gain = beta / alpha(n);      // Kalman gain for this scalar measurement
  K_.col(meas) = k_gain;

  Real residual = dz_(meas);
  if (meas > 0) residual -= H_.row(meas).dot(dx_);   // correct for prior updates
  dx_ += k_gain * residual.val();

  D_diag_ = D_plus;
  U_ = U_plus;
}
x_ = x_ + dx_;
\end{lstlisting}

\begin{lupntnote}
\cpp{UDUEKF} inherits \cpp{EKF::RemoveOutliers}, so the gate of
\cref{eq:flt-outlier} runs on the \emph{dense} innovation covariance
$\mat{S} = \mat{R} + \mat{H}\mat{P}\mat{H}\transpose$ formed once before the
Carlson loop. Only the covariance recursion is factorized; the screening is
not.
\end{lupntnote}

\section{Square-Root Information Filter}
\label{sec:flt-srif}

\cpp{SRIF} is a subclass of \cpp{EKF} that stores the \emph{information} in
square-root form. Let $\mat{R}_{\text{info}}$ be upper triangular with
\begin{equation}
  \label{eq:flt-srif-def}
  \mat{R}_{\text{info}}\transpose\,\mat{R}_{\text{info}} = \mat{P}^{-1},
  \qquad\text{equivalently}\qquad
  \mat{P} = \mat{R}_{\text{info}}^{-1}\,\mat{R}_{\text{info}}^{-\mathsf{T}} .
\end{equation}
The factor is initialized from the user's covariance by a Cholesky
factorization of $\mat{P}^{-1}$ in \cpp{SRIF::SetInfoSqrt}. Like the UD filter,
the SRIF halves the condition number of the stored object; unlike the UD
filter, it handles a full (non-diagonal) $\mat{R}$ natively and admits a
genuinely infinite prior variance --- a zero row of $\mat{R}_{\text{info}}$ ---
which is why it is the natural choice for a cold start with no useful a-priori
state. The dense $\mat{P}$ and $\hat{\vec{x}}$ are re-synchronized after every
step so that the inherited outlier gate, fault detection, and RTS smoother work
unchanged \citep[Ch.~5]{tapley2004}.

\subsection{Time update (Dyer--McReynolds)}

Let $\mat{S}$ be a square root of the process noise, $\mat{Q} = \mat{S}
\mat{S}\transpose$ (a Cholesky factor; $\mat{S} = \mat{0}$ when there is no
process noise), and let $\mat{R}_\Phi = \mat{R}_{\text{info}}\,\mat{\Phi}^{-1}$
with $\mat{\Phi} = \mat{F}$ the STM. A Householder QR triangularization of the
$2n \times 2n$ array
\begin{equation}
  \label{eq:flt-srif-time}
  \begin{bmatrix}
    \mat{I} & \mat{0} \\
    -\mat{R}_\Phi \mat{S} & \mat{R}_\Phi
  \end{bmatrix}
  \;\xrightarrow{\ \text{QR}\ }\;
  \begin{bmatrix}
    * & * \\
    \mat{0} & \bar{\mat{R}}_{\text{info}}
  \end{bmatrix}
\end{equation}
zeros the bottom-left block; the bottom-right block is the a-priori information
square root at $t_f$. The upper blocks marked $*$ are the (discarded)
information about the process-noise realization itself. Because the estimate is
re-centered on $\hat{\vec{x}}$ at every step, the data column is identically
zero and is omitted from the array.

\begin{lstlisting}[language=cpp, caption={SRIF time update.},
                   label={lst:flt-srif-predict}]
// cpp/lupnt/numerics/filters/srif.cc :: SRIF::Predict
MatXd S = MatXd::Zero(n, n);
if (!Q_.isZero()) { Eigen::LLT<MatXd> llt(Q_); S = llt.matrixL(); }  // Q = S S^T

MatXd RPhiInv = R_info_ * F_.inverse();
MatXd A = MatXd::Zero(2 * n, 2 * n);
A.topLeftCorner(n, n)     = MatXd::Identity(n, n);
A.bottomLeftCorner(n, n)  = -RPhiInv * S;
A.bottomRightCorner(n, n) = RPhiInv;
R_info_ = QrRFactor(A).bottomRightCorner(n, n);

MatXd R_inv = R_info_.inverse();
P_ = R_inv * R_inv.transpose();
\end{lstlisting}

\subsection{Measurement update (Householder information update)}

The measurement rows are first whitened. With the Cholesky factorization
$\mat{R} = \mat{L}\mat{L}\transpose$, the data equation
$\delta \vec{z} = \mat{H}\,\delta\vec{x} + \vec{v}$,
$\vec{v} \sim \mathcal{N}(\vec{0}, \mat{R})$ becomes a unit-variance system,
\begin{equation}
  \label{eq:flt-srif-whiten}
  \mat{H}_w = \mat{L}^{-1}\mat{H},
  \qquad
  \delta \vec{z}_w = \mat{L}^{-1}\,\delta \vec{z},
  \qquad
  \delta \vec{z}_w = \mat{H}_w\,\delta \vec{x} + \vec{v}_w,
  \quad \vec{v}_w \sim \mathcal{N}(\vec{0}, \mat{I}) .
\end{equation}
Both solves are triangular. The whitened rows are then folded into the
information array by a second Householder QR,
\begin{equation}
  \label{eq:flt-srif-meas}
  \begin{bmatrix}
    \mat{R}_{\text{info}} & \mat{0} \\
    \mat{H}_w & \delta \vec{z}_w
  \end{bmatrix}
  \;\xrightarrow{\ \text{QR}\ }\;
  \begin{bmatrix}
    \hat{\mat{R}}_{\text{info}} & \hat{\vec{b}} \\
    \mat{0} & *
  \end{bmatrix} ,
\end{equation}
after which the state correction is a single back-substitution against the
updated triangular factor,
\begin{equation}
  \label{eq:flt-srif-dx}
  \delta \vec{x} = \hat{\mat{R}}_{\text{info}}^{-1}\,\hat{\vec{b}},
  \qquad
  \hat{\vec{x}}^{+} = \hat{\vec{x}}^{-} + \delta \vec{x} .
\end{equation}
The orthogonal factor $\mat{Q}$ of the QR is never assembled: only the
triangular $\mat{R}$ factor of the augmented array is needed, and the
right-hand column is triangularized along with it so that its top rows emerge
as $\mat{Q}\transpose \delta \vec{z}_w$. The residual entry marked $*$ in
\cref{eq:flt-srif-meas} is the square root of the post-fit sum of squares, a
free goodness-of-fit statistic. No gain matrix and no matrix inverse of a
measurement-sized object appear anywhere in the update; the whole step is
backward stable.

\begin{lstlisting}[language=cpp, caption={SRIF Householder information update.},
                   label={lst:flt-srif-update}]
// cpp/lupnt/numerics/filters/srif.cc :: SRIF::InformationUpdate
Eigen::LLT<MatXd> llt(R_);
MatXd L = llt.matrixL();                                            // R = L L^T
MatXd H_w  = L.template triangularView<Eigen::Lower>().solve(H_);   // m x n
VecXd dz_w = L.template triangularView<Eigen::Lower>().solve(dz_);  // m

MatXd A(n + m, n + 1);
A.topLeftCorner(n, n)    = R_info_;
A.block(0, n, n, 1).setZero();
A.bottomLeftCorner(m, n) = H_w;
A.block(n, n, m, 1)      = dz_w;

MatXd Rtri = QrRFactor(A);
R_info_ = Rtri.topLeftCorner(n, n);
VecXd b_hat = Rtri.block(0, n, n, 1);

dx_ = R_info_.template triangularView<Eigen::Upper>().solve(b_hat);
x_  = x_ + dx_;

MatXd R_inv = R_info_.inverse();
P_ = R_inv * R_inv.transpose();
\end{lstlisting}

\begin{lupntwarning}
\Cref{eq:flt-srif-time} inverts the STM. This is exact for the invertible flows
of orbital dynamics but is a real cost ($O(n^3)$ per step) and a real failure
mode for a near-singular $\mat{F}$, for example one containing an exactly
non-observable direction with zero eigenvalue. The UD filter's time update has
no such requirement.
\end{lupntwarning}

\section{Fixed-Interval Smoothing}
\label{sec:flt-smoother}

For post-processing, \cpp{EKF} logs the prior and posterior states,
covariances, and STMs at every epoch (\cpp{InitializeLogger},
\cpp{LogFilterEstimate}) and runs a backward Rauch--Tung--Striebel (RTS)
recursion over the log. With $N$ the index of the last epoch, the recursion is
initialized at $\hat{\vec{x}}_{N|N} = \hat{\vec{x}}^{+}_N$,
$\mat{P}_{N|N} = \mat{P}^{+}_N$ and steps backwards with the gain
\begin{align}
  \label{eq:flt-rts-gain}
  \mat{A}_k &= \mat{P}_{k|k}\,\mat{F}_{k+1}\transpose\,\mat{P}_{k+1|k}^{-1}, \\
  \label{eq:flt-rts-x}
  \hat{\vec{x}}_{k|N} &= \hat{\vec{x}}_{k|k}
    + \mat{A}_k\bigl(\hat{\vec{x}}_{k+1|N} - \hat{\vec{x}}_{k+1|k}\bigr), \\
  \label{eq:flt-rts-p}
  \mat{P}_{k|N} &= \mat{P}_{k|k}
    + \mat{A}_k\bigl(\mat{P}_{k+1|N} - \mat{P}_{k+1|k}\bigr)\mat{A}_k\transpose .
\end{align}
Here $\mat{F}_{k+1}$ is the STM of the propagation from $t_k$ to $t_{k+1}$,
stored at index $k+1$ of the log. \Cref{eq:flt-rts-p} is a correction of
definite sign only in the sense that
$\mat{P}_{k+1|N} \preceq \mat{P}_{k+1|k}$: the smoothed covariance is never
larger than the filtered one.

\begin{lstlisting}[language=cpp, caption={RTS backward recursion.},
                   label={lst:flt-smoother}]
// cpp/lupnt/numerics/filters/ekf.cc :: EKF::UpdateSmoother
int k = tidx;
MatXd P_kk  = P_pos_log_[k];
MatXd P_k1k = P_prior_log_[k + 1];
MatXd stm_k = stm_log_[k + 1];
VecXd x_kk  = x_pos_log_[k];
VecXd x_k1k = x_prior_log_[k + 1];

MatXd A_k = P_kk * stm_k.transpose() * P_k1k.inverse();

x_sm_[k] = x_kk + A_k * (x_sm_[k + 1] - x_k1k);
P_sm_[k] = P_kk + A_k * (P_sm_[k + 1] - P_k1k) * A_k.transpose();
\end{lstlisting}

\Cref{eq:flt-rts-gain,eq:flt-rts-x,eq:flt-rts-p} are exact only when the
measurements at each epoch depend on the state at that epoch alone. That
assumption fails for time-differenced carrier phase, and \cref{sec:flt-cloning}
gives the replacement.

\section{Stochastic-Cloning UD Filter}
\label{sec:flt-cloning}

Time-differenced carrier-phase (TDCP) measurements depend on the state at
\emph{two} epochs. Differencing the carrier phase between consecutive epochs
cancels the (unknown, constant) integer ambiguity and leaves a precise
delta-range observable, at the cost of an observation model
\begin{equation}
  \label{eq:flt-sc-tdcp}
  \vec{z}_{k} = h_1(\vec{x}_{k}) - h_2(\vec{x}_{k-1}) + \vec{v}_{k}
\end{equation}
that is no longer a function of the current state alone
\citep{iiyama2024navi,iiyama2026scud}. Treating \cref{eq:flt-sc-tdcp} as if it
were a current-state measurement --- by folding the previous epoch's
contribution into the noise --- both mis-states the noise covariance and
destroys the correlation structure the differencing created. The remedy is
stochastic cloning \citep{roumeliotis2002cloning}: carry a copy of the previous
state inside the filter so that the delayed-state dependency becomes an
ordinary linear dependency on an augmented state.

\cpp{UDUStochasticCloningEKF} implements this while keeping the full
$2n \times 2n$ augmented covariance in UD form.

\subsection{Augmented state and factors}

The augmented state and its UD-factored, time-propagated covariance are
\citep[Ch.~6.2.1, Eqs.~6.18--6.23]{iiyama2026thesis}
\begin{equation}
  \label{eq:flt-sc-state}
  \vec{x}^{\mathrm{sc}}_{k-1}
  = \begin{bmatrix}
      \hat{\vec{x}}_{k-1|k-1} \\ \hat{\vec{x}}_{k-1|k-1}
    \end{bmatrix},
\end{equation}
\begin{equation}
  \label{eq:flt-sc-ud}
  \hat{\mat{U}}^{\mathrm{sc}}_{k|k-1}
  = \begin{bmatrix}
      \mat{G}_k & \mat{\Phi}_{k,k-1}\hat{\mat{U}}_{k-1|k-1} \\
      \mat{0}   & \hat{\mat{U}}_{k-1|k-1}
    \end{bmatrix},
  \qquad
  \hat{\mat{D}}^{\mathrm{sc}}_{k|k-1}
  = \begin{bmatrix}
      \mat{Q}^{D}_k & \mat{0} \\
      \mat{0}       & \hat{\mat{D}}_{k-1|k-1}
    \end{bmatrix} ,
\end{equation}
so that the top block is the propagated state, transported by the variational
STM $\mat{\Phi}_{k,k-1}$, and the bottom block is the frozen clone. In \lupnt{}
the augmented STM is assembled explicitly as
\begin{equation}
  \label{eq:flt-sc-stm}
  \mat{F} =
  \begin{bmatrix}
    \mat{F}_{\text{base}} & \mat{0} \\
    \mat{I}               & \mat{0}
  \end{bmatrix}
  \in \mathbb{R}^{2n \times 2n} ,
\end{equation}
whose bottom row copies the pre-propagation state into the clone slot, and the
Thornton MWGS of \cref{eq:flt-udu-mwgs} is applied to the augmented system
with the process noise acting on the current block only:
\begin{equation}
  \label{eq:flt-sc-mwgs}
  \mat{E} = \begin{bmatrix} \mat{G} \\ \mat{0} \end{bmatrix}
    \in \mathbb{R}^{2n \times m},
  \qquad
  \mat{Y} = \begin{bmatrix} \mat{F}\mat{U} & \mat{E} \end{bmatrix},
  \qquad
  \mat{Y}\tilde{\mat{D}}\mat{Y}\transpose
  = \mat{F}\mat{P}\mat{F}\transpose
    + \operatorname{blkdiag}\bigl(\mat{G}\mat{Q}\mat{G}\transpose, \mat{0}\bigr) .
\end{equation}
The clone block therefore receives no process noise --- correctly, since it is
a stored value, not a propagated one --- and the factorization is still
computed without ever forming the $2n \times 2n$ covariance.

\begin{lstlisting}[language=cpp, caption={Augmented STM assembly and MWGS time update.},
                   label={lst:flt-sc-predict}]
// cpp/lupnt/numerics/filters/udu_filter.cc :: UDUStochasticCloningEKF::Predict
State x_prev_current = x_.head(n);
MatXd F_base;
State x_current = f_dyn_(x_prev_current, t_, t, u, &F_base);

State x_pred(2 * n);
x_pred.head(n) = x_current;        // propagated block
x_pred.tail(n) = x_prev_current;   // cloned block
x_ = x_pred;

const int n2 = 2 * n;
F_.setZero(n2, n2);
F_.topLeftCorner(n, n)    = F_base;                 // [ F_base 0 ]
F_.bottomLeftCorner(n, n) = MatXd::Identity(n, n);  // [ I      0 ]

MatXd E = MatXd::Zero(n2, m);
E.topRows(n) = G_;                 // process noise on the current block only
MatXd Y = MatXd::Zero(n2, n2 + m);
Y.leftCols(n2) = F_ * U_;
Y.rightCols(m) = E;
MatXd D_tilde = MatXd::Zero(n2 + m, n2 + m);
D_tilde.topLeftCorner(n2, n2)   = D_diag_.asDiagonal();
D_tilde.bottomRightCorner(m, m) = Q_;
std::tie(D_diag_, U_) = ModifiedGramSchmidt(D_tilde, Y);
\end{lstlisting}

\subsection{Delayed-state Jacobian and update ordering}

The measurement Jacobian with respect to the augmented state splits into
current- and cloned-state parts
\citep[Ch.~6.2.3, Eq.~6.25]{iiyama2026thesis},
\begin{equation}
  \label{eq:flt-sc-jac}
  \mat{H}_{k,m}
  = \begin{bmatrix} \mat{H}^{(1)}_{k,m} & \mat{H}^{(2)}_{k,m} \end{bmatrix},
  \qquad
  \mat{H}^{(1)}_{k,m} = \frac{\partial h}{\partial \vec{x}_{k}},
  \qquad
  \mat{H}^{(2)}_{k,m} = \frac{\partial h}{\partial \vec{x}_{k-1}} ,
\end{equation}
with $\mat{H}^{(2)} = \mat{0}$ for any current-state-only observable such as a
pseudorange or a Doppler shift. For the TDCP model of \cref{eq:flt-sc-tdcp},
$\mat{H}^{(1)} = \partial h_1/\partial \vec{x}_k$ and
$\mat{H}^{(2)} = -\partial h_2/\partial \vec{x}_{k-1}$. Both blocks come from
the same automatic-differentiation pass over the augmented state, so no
delayed-state model needs to be differentiated by hand.

Measurements are assimilated by the same Carlson rank-one recursion of
\cref{eq:flt-udu-carlson-f,eq:flt-udu-alpha,eq:flt-udu-d,eq:flt-udu-u}, applied
to the $2n$-dimensional factors
\citep[Eqs.~6.24--6.30]{iiyama2026thesis}. Ordering matters for cost, not for
the result: because the current-state-only rows have
$\mat{H}^{(2)} = \mat{0}$, their $\vec{f} = \mat{U}\transpose\vec{h}$ has
support only in the leading $n$ entries, so processing the delayed-state (TDCP)
rows \emph{first} lets every subsequent current-state update skip the cloned
sub-blocks $\hat{\mat{U}}_{22}$ entirely
\citep[Ch.~6.2.4]{iiyama2026thesis}.

\subsection{The delayed-state smoother}

The standard RTS recursion of \cref{eq:flt-rts-gain,eq:flt-rts-x} rests on the
Markov assumption \citep[Eq.~6.36]{iiyama2026thesis}, which
\cref{eq:flt-sc-tdcp} violates: the measurement at epoch $k+1$ is correlated
with the state at epoch $k$ \citep[Eq.~6.37]{iiyama2026thesis}, so
$\hat{\vec{x}}_{k|k}$ is not a sufficient statistic for the past. The thesis
therefore derives the smoother gain from the joint stochastic-cloning
posterior \citep[Eqs.~6.38--6.41]{iiyama2026thesis}, giving
\begin{equation}
  \label{eq:flt-sc-gain}
  \hat{\vec{x}}_{k|N} = \hat{\vec{x}}_{k+1|k+1}
    + \mat{J}_k\bigl(\hat{\vec{x}}_{k+1|N} - \hat{\vec{x}}_{k+1|k+1}\bigr),
  \qquad
  \mat{J}_k = \hat{\mat{P}}_{k+1,k|k+1}\transpose\,
              \hat{\mat{P}}_{k+1|k+1}^{-1} ,
\end{equation}
where $\hat{\mat{P}}_{k+1,k|k+1}$ is the \emph{posterior} cross-covariance
between the current and cloned blocks --- that is, the off-diagonal block of
the augmented covariance after the update. This is exactly the
measurement-induced correlation that $\mat{A}_k$ in \cref{eq:flt-rts-gain}
misses, since $\mat{A}_k$ reconstructs the cross-covariance from the STM alone.
Both quantities in \cref{eq:flt-sc-gain} are available directly from the stored
augmented UD factors, and the gain is obtained without forming an explicit
inverse \citep[Eq.~6.42]{iiyama2026thesis} by the three-stage solve
\begin{equation}
  \label{eq:flt-sc-solve}
  \mat{U}\diag(\vec{D})\mat{U}\transpose \mat{X} = \mat{B}
  \quad\Longrightarrow\quad
  \underbrace{\mat{U}\mat{Y} = \mat{B}}_{\text{back substitution}},
  \quad
  \underbrace{\diag(\vec{D})\mat{Z} = \mat{Y}}_{\text{diagonal scaling}},
  \quad
  \underbrace{\mat{U}\transpose\mat{X} = \mat{Z}}_{\text{forward substitution}} ,
\end{equation}
implemented as \cpp{UDUSolve} in
\file{cpp/lupnt/numerics/filters/udu_utils.cc}. Both triangular systems are
unit-triangular, so \cref{eq:flt-sc-solve} costs $O(n^2)$ per right-hand side
and touches no division except by $\vec{D}$.

\begin{lupntnote}
\cpp{EKF::UpdateSmoother} implements the \emph{standard} RTS recursion
(\cref{eq:flt-rts-gain}), which is exact when measurements depend only on the
current state. The cross-covariance gain $\mat{J}_k$ of
\cref{eq:flt-sc-gain} is the thesis extension required when delayed-state
(TDCP) measurements are present; it operates on the augmented
stochastic-cloning covariance maintained by \cpp{UDUStochasticCloningEKF}.
Applying the standard gain to a TDCP run does not merely lose optimality --- it
produces an inconsistent covariance.
\end{lupntnote}

\section{Batch Weighted Least Squares}
\label{sec:flt-batch}

\cpp{RunBatchFilter} (\file{cpp/lupnt/numerics/filters/batch_filter.cc}) solves
the nonlinear weighted least-squares problem by Gauss--Newton iteration. There
is no time propagation: the whole measurement set is treated as a function of a
single epoch state, which makes it the right tool for snapshot positioning and
for initializing a recursive filter.

At each iteration all residuals, Jacobians, and weights are stacked,
\begin{equation}
  \label{eq:flt-batch-stack}
  \delta \vec{z} = \vec{z} - h(\hat{\vec{x}}),
  \qquad
  \mat{H} = \frac{\partial h}{\partial \vec{x}}\bigg|_{\hat{\vec{x}}},
  \qquad
  \mat{W} = \diag\bigl(1/\sigma_i^2\bigr) ,
\end{equation}
and the normal equations are solved for the correction,
\begin{equation}
  \label{eq:flt-batch-normal}
  \bigl(\mat{H}\transpose \mat{W} \mat{H}\bigr)\,\delta \vec{x}
    = \mat{H}\transpose \mat{W}\,\delta \vec{z},
  \qquad
  \hat{\vec{x}} \leftarrow \hat{\vec{x}} + \delta \vec{x} ,
\end{equation}
by an LDLT factorization (\cpp{SolveWeightedLeastSquares}). When
initialization constraints are enabled (\code{use\_initialization}), a soft
prior is appended by augmenting the system,
\begin{equation}
  \label{eq:flt-batch-prior}
  \mat{H} \leftarrow \begin{bmatrix} \mat{H} \\ \mat{I} \end{bmatrix},
  \qquad
  \delta \vec{z} \leftarrow
    \begin{bmatrix} \delta \vec{z} \\ \vec{x}_0 - \hat{\vec{x}} \end{bmatrix},
  \qquad
  \mat{W} \leftarrow
    \operatorname{blkdiag}\bigl(\mat{W},\ \diag(1/\sigma_{0}^2)\bigr) ,
\end{equation}
which is the exact equivalent of a Gaussian prior $\mathcal{N}(\vec{x}_0,
\diag(\sigma_0^2))$ and regularizes an otherwise rank-deficient geometry.
Iteration stops when $\lVert \delta \vec{x}\rVert$ falls below
\code{convergence\_tol} or the iteration cap \code{max\_iterations} is reached;
each iteration records its correction norm and weighted residual RMS in
\cpp{iteration\_history}.

The reported covariance is the inverse of the \emph{measurement} information
only --- the prior of \cref{eq:flt-batch-prior} is deliberately excluded, so
that the number reflects the observation geometry rather than the strength of
the initialization --- and is computed by an SVD pseudo-inverse so that a
rank-deficient geometry returns a finite answer instead of failing:
\begin{equation}
  \label{eq:flt-batch-cov}
  \mat{P} = \bigl(\mat{H}\transpose \mat{W} \mat{H}\bigr)^{+},
  \qquad
  \text{PDOP} = \sqrt{\Tr\bigl(\mat{P}_{[0:3,\,0:3]}\bigr)} .
\end{equation}
The PDOP of \cref{eq:flt-batch-cov} assumes that the first three state
components are position and that $\mat{W}$ carries the measurement variances,
so it is a metre-valued position-error standard deviation rather than a
dimensionless dilution factor.

\begin{lstlisting}[language=cpp, caption={Batch normal-equation solve and covariance.},
                   label={lst:flt-batch}]
// cpp/lupnt/numerics/filters/batch_filter.cc :: SolveWeightedLeastSquares
MatXd W    = weights.asDiagonal();
MatXd HTW  = H.transpose() * W;
return (HTW * H).ldlt().solve(HTW * residuals);   // (H^T W H) dx = H^T W r

// cpp/lupnt/numerics/filters/batch_filter.cc :: RunBatchFilter
// final covariance -- SVD pseudo-inverse of the measurement information
Eigen::JacobiSVD<MatXd> svd(information_matrix,
                            Eigen::ComputeFullU | Eigen::ComputeFullV);
covariance_matrix = svd.solve(MatXd::Identity(n_state, n_state));

// cpp/lupnt/numerics/filters/batch_filter.cc :: CalculatePDOP
Mat3d pos_cov = covariance_matrix.block<3, 3>(0, 0);
return std::sqrt(pos_cov.trace());
\end{lstlisting}

\section{Process-Noise Models}
\label{sec:flt-noise}

The process noise is the filter's statement about what the dynamics model does
\emph{not} capture. \lupnt{} ships closed-form $\mat{Q}$ blocks for the two
dominant cases --- unmodelled acceleration on the orbit and oscillator noise on
the clock --- plus a sliding-window adaptive scheme.

\subsection{Continuous white-noise acceleration}
\label{sec:flt-cwna}

Model the unmodelled acceleration as a zero-mean, white, continuous process
with per-axis power spectral density $q_a$ acting on the velocity,
\begin{equation}
  \label{eq:flt-cwna-sde}
  \dd \begin{bmatrix} \vec{r} \\ \vec{v} \end{bmatrix}
  = \begin{bmatrix} \vec{v} \\ \vec{0} \end{bmatrix} \dd t
  + \begin{bmatrix} \mat{0} \\ \mat{I}_3 \end{bmatrix} \sqrt{q_a}\,\dd \vec{w},
  \qquad
  \mathbb{E}\bigl[\dd \vec{w}\,\dd \vec{w}\transpose\bigr] = \mat{I}_3\,\dd t .
\end{equation}
Integrating the state-transition form of \cref{eq:flt-cwna-sde} over a step
$\Delta t$ against the double-integrator STM gives the standard $6\times 6$
block for a Cartesian state $[\vec{r}\transpose\ \vec{v}\transpose]\transpose$,
\begin{equation}
  \label{eq:flt-cwna}
  \mat{Q} = q_a
  \begin{bmatrix}
    \tfrac{\Delta t^3}{3}\,\mat{I}_3 & \tfrac{\Delta t^2}{2}\,\mat{I}_3 \\[6pt]
    \tfrac{\Delta t^2}{2}\,\mat{I}_3 & \Delta t\,\mat{I}_3
  \end{bmatrix},
  \qquad
  [q_a] = \si{\meter\squared\per\second\cubed} ,
\end{equation}
returned by \cpp{CwnaProcessNoise(dt, accel\_psd)}. The block is
\emph{singular in the noise directions}: it has rank 3, not 6, since a single
acceleration process drives both the position and velocity blocks. That is why
\cref{eq:flt-cwna} may not be handed to the UD or stochastic-cloning filters as
a diagonal $\mat{Q}$; the correct route there is
$\mat{Q} = q_a \mat{I}_3$ with the mapping matrix
$\mat{G} = [\tfrac{\Delta t^2}{2}\mat{I}_3;\ \Delta t\,\mat{I}_3]$ of
\cref{eq:flt-gmap}.

\begin{lstlisting}[language=cpp, caption={Continuous white-noise acceleration block.},
                   label={lst:flt-cwna}]
// cpp/lupnt/numerics/filters/srif.cc :: CwnaProcessNoise
MatXd Q = MatXd::Zero(6, 6);
const double q3 = accel_psd * dt * dt * dt / 3.0;
const double q2 = accel_psd * dt * dt / 2.0;
const double q1 = accel_psd * dt;
for (int k = 0; k < 3; ++k) {
  Q(k, k)         = q3;
  Q(k, k + 3)     = q2;
  Q(k + 3, k)     = q2;
  Q(k + 3, k + 3) = q1;
}
return Q;
\end{lstlisting}

\subsection{Clock process noise}
\label{sec:flt-clock-noise}

\cpp{ClockDynamics} (\file{cpp/lupnt/dynamics/clock_dynamics.cc}) propagates a
two- or three-state clock error with the polynomial transition matrices
\begin{equation}
  \label{eq:flt-clock-phi}
  \mat{\Phi}_2(\Delta t) = \begin{bmatrix} 1 & \Delta t \\ 0 & 1 \end{bmatrix},
  \qquad
  \mat{\Phi}_3(\Delta t) = \begin{bmatrix}
    1 & \Delta t & \tfrac{\Delta t^2}{2} \\
    0 & 1        & \Delta t \\
    0 & 0        & 1
  \end{bmatrix} ,
\end{equation}
acting on the state $(b, \dot{b})$ or $(b, \dot{b}, \ddot{b})$ --- clock bias,
fractional frequency offset, and frequency drift rate.

The noise model is the standard two- or three-state clock diffusion of
\citet{zucca2005clock}: three independent Wiener processes with diffusion
coefficients $(q_1, q_2, q_3)$ drive the phase, the frequency, and the
frequency drift respectively,
\begin{equation}
  \label{eq:flt-clock-sde}
  \dd \begin{bmatrix} b \\ \dot{b} \\ \ddot{b} \end{bmatrix}
  = \begin{bmatrix} \dot{b} \\ \ddot{b} \\ 0 \end{bmatrix} \dd t
  + \begin{bmatrix}
      \sqrt{q_1} & 0 & 0 \\
      0 & \sqrt{q_2} & 0 \\
      0 & 0 & \sqrt{q_3}
    \end{bmatrix} \dd \vec{w} .
\end{equation}
Integrating $\mat{Q} = \int_0^{\Delta t} \mat{\Phi}(\Delta t - s)\,
\mat{B}\mat{B}\transpose \mat{\Phi}(\Delta t - s)\transpose \dd s$ in closed
form gives, for the two-state model (\cpp{TwoStateNoise}),
\begin{equation}
  \label{eq:flt-clock-q2}
  \mat{Q}_2 =
  \begin{bmatrix}
    q_1 \Delta t + q_2\dfrac{\Delta t^3}{3} & q_2\dfrac{\Delta t^2}{2} \\[10pt]
    q_2\dfrac{\Delta t^2}{2}                & q_2 \Delta t
  \end{bmatrix},
\end{equation}
and, for the three-state model (\cpp{ThreeStateNoise}),
\begin{equation}
  \label{eq:flt-clock-q3}
  \mat{Q}_3 =
  \begin{bmatrix}
    q_1\Delta t + q_2\dfrac{\Delta t^3}{3} + q_3\dfrac{\Delta t^5}{20}
      & q_2\dfrac{\Delta t^2}{2} + q_3\dfrac{\Delta t^4}{8}
      & q_3\dfrac{\Delta t^3}{6} \\[10pt]
    q_2\dfrac{\Delta t^2}{2} + q_3\dfrac{\Delta t^4}{8}
      & q_2\Delta t + q_3\dfrac{\Delta t^3}{3}
      & q_3\dfrac{\Delta t^2}{2} \\[10pt]
    q_3\dfrac{\Delta t^3}{6}
      & q_3\dfrac{\Delta t^2}{2}
      & q_3\Delta t
  \end{bmatrix} .
\end{equation}
\Cref{eq:flt-clock-q2} is the leading $2\times 2$ block of
\cref{eq:flt-clock-q3} with $q_3 = 0$, as it must be. The dimensions are
$[q_1] = \si{\second}$, $[q_2] = \si{\per\second}$, and
$[q_3] = \si{\per\second\cubed}$, so that every entry of
\cref{eq:flt-clock-q3} carries the square of the corresponding state's unit.

The diffusion coefficients relate directly to the Allan variance of the
oscillator, which is how they are measured in practice
\citep{zucca2005clock}:
\begin{equation}
  \label{eq:flt-clock-allan}
  \sigma_y^2(\tau) = \frac{q_1}{\tau} + \frac{q_2\,\tau}{3}
                   + \frac{q_3\,\tau^3}{20} ,
\end{equation}
whose three terms are the white-frequency ($\tau^{-1}$), random-walk-frequency
($\tau^{+1}$), and random-run ($\tau^{+3}$) branches of the classical Allan
deviation plot. \Cref{tab:flt-clock} lists the coefficients hard-coded in
\cpp{ClockDynamics::GetClockValues}.

\begin{table}[H]
  \centering
  \caption{Clock diffusion coefficients in
           \file{cpp/lupnt/dynamics/clock_dynamics.cc}
           (\cpp{ClockDynamics::GetClockValues}).}
  \label{tab:flt-clock}
  \begin{tabularx}{\textwidth}{lllLL}
    \toprule
    \cpp{ClockModel} & $q_1$ [\si{\second}] & $q_2$ [\si{\per\second}]
      & $q_3$ [\si{\per\second\cubed}] & Source \\
    \midrule
    \code{OCXO}      & \num{3.881e-25} & \num{4.052e-28} & \num{4.917e-55}
      & Oven-controlled crystal; first LCRNS satellite design study \\
    \code{USO}       & \num{1.2e-22}   & \num{1.58e-26}  & \num{5.9e-75}
      & Ultra-stable oscillator, as flown on MMS \\
    \code{CSAC}      & \num{5.8e-22}   & \num{1.5e-29}   & \num{1e-45}
      & Chip-scale atomic clock \\
    \code{MINI\_RAFS}& \num{1.0e-22}   & \num{1.21e-30}  & \num{4.917e-55}
      & Miniaturized rubidium atomic frequency standard \\
    \code{RAFS}      & \num{3.70e-24}  & \num{1.87e-33}  & \num{7.56e-59}
      & Rubidium AFS; NASA Gateway autonomous-navigation study \\
    \code{DSAC}      & \num{4.23e-26}  & \num{6.19e-38}  & \num{2.24e-61}
      & Deep Space Atomic Clock \\
    \bottomrule
  \end{tabularx}
\end{table}

\begin{lupntnote}
The clock bias may be carried in seconds, metres, or kilometres, selected by
\cpp{ClockBiasUnit}. \Cref{eq:flt-clock-q2,eq:flt-clock-q3} are derived in
seconds and rescaled by the overloads that take a unit,
\begin{equation}
  \label{eq:flt-clock-scale}
  \mat{Q} \leftarrow \sigma^2 \mat{Q},
  \qquad
  \sigma = \texttt{SecondsToBiasUnitScale} =
  \begin{cases}
    1 & \text{(seconds)}, \\
    c & \text{(metres)}, \\
    c \cdot 10^{-3} & \text{(kilometres)},
  \end{cases}
\end{equation}
with $c$ the speed of light. All three clock states scale by the same factor,
so a single scalar $\sigma^2$ suffices for the whole block. Mixing a
metre-valued clock state with a second-valued $\mat{Q}$ is silent and
catastrophic --- always obtain $\mat{Q}$ from the unit-aware overload.
\end{lupntnote}

This is the clock half of the coupled orbit--clock process noise of
\citep[Ch.~6.1.2, Eq.~6.6]{iiyama2026thesis}; the orbit half is
\cref{eq:flt-cwna}. For a coupled state the full block is
$\operatorname{blkdiag}(\mat{Q}_{\text{orbit}}, \mat{Q}_{\text{clock}})$, which
is not diagonal, so a UD filter must receive it through the mapping matrix
$\mat{G}$.

\begin{lstlisting}[language=cpp, caption={Three-state clock process noise and the unit rescaling.},
                   label={lst:flt-clock}]
// cpp/lupnt/dynamics/clock_dynamics.cc :: ClockDynamics::ThreeStateNoise
Q(0, 0) = q1 * dt + q2 * dt3 / 3 + q3 * dt5 / 20.0;
Q(0, 1) = q2 * dt2 / 2.0 + q3 * dt4 / 8.0;
Q(0, 2) = q3 * dt3 / 6.0;
Q(1, 0) = Q(0, 1);
Q(1, 1) = q2 * dt + q3 * dt3 / 3.0;
Q(1, 2) = q3 * dt2 / 2.0;
Q(2, 0) = Q(0, 2);
Q(2, 1) = Q(1, 2);
Q(2, 2) = q3 * dt;

// bias-unit overload: Q_unit = sigma^2 Q, sigma = SecondsToBiasUnitScale(unit)
Real scale = SecondsToBiasUnitScale(unit);
return ThreeStateNoise(clk_model, dt) * scale * scale;
\end{lstlisting}

\subsection{Adaptive process noise}
\label{sec:flt-adaptive}

The class \cpp{ProcessNoise}
(\file{cpp/lupnt/numerics/filters/adaptive_process_noise.cc}) re-estimates the
acceleration noise level online from the filter's own recent behaviour, under
one of three algorithms selected by \cpp{ProcessNoiseAlgorithm}:
\code{SNC} (state noise compensation, i.e.\ a fixed $\mat{Q}_a$ and no
adaptation), \code{ASNC} (adaptive SNC: a piecewise-constant white-noise
acceleration model on the position/velocity states, whose coefficient basis is
shared across the three axes), or \code{ADMC} (adaptive dynamic model
compensation: a first-order Gauss--Markov acceleration state of time constant
$1/\beta$, with a per-axis coefficient basis). Both adaptive variants solve a
per-axis diagonal $\mat{Q}_a$; they differ in the acceleration model and hence
in whether that coefficient basis is shared or per-axis.

Each step feeds the tuple $(\delta \vec{x}, \mat{\Sigma}_{\delta x},
\bar{\mat{P}}, \mat{P}^{+})$ into \cpp{ProcessNoise::Update}, which maintains a
sliding window of length $N_w$ after an initial burn-in of $N_{\text{wait}}$
steps. The window yields the empirical process-noise estimate
\begin{equation}
  \label{eq:flt-adaptive-qhat}
  \hat{\mat{Q}} = \frac{1}{N_w}\sum_{i=1}^{N_w}
    \Bigl(\mat{P}^{+}_i - \bar{\mat{P}}_i
          + \delta \vec{x}_i\,\delta \vec{x}_i\transpose\Bigr) ,
\end{equation}
which is the sample second moment of the corrections, debiased by the
covariance the filter already claims for them. \Cref{eq:flt-adaptive-qhat} is
then matched, in a weighted least-squares sense whose weights come from the
fourth-moment matrix accumulated alongside it, against the structural
coefficients of \cref{eq:flt-cwna}; solving for the per-axis acceleration PSD
gives $\mat{Q}_a$, which is clamped to the configured limits,
\begin{equation}
  \label{eq:flt-adaptive-clamp}
  \bigl[\mat{Q}_a\bigr]_{ii} \leftarrow
    \min\Bigl(\max\bigl(\bigl[\mat{Q}_a\bigr]_{ii},\ \sigma^2_{\min}\bigr),\
              \sigma^2_{\max}\Bigr) ,
\end{equation}
and optionally smoothed by an exponential filter with rate $\alpha$,
$\mat{Q}_a \leftarrow (1-\alpha)\mat{Q}_a + \alpha \hat{\mat{Q}}_a$. Any window
entry with a negative variance on the diagonal resets the accumulator rather
than poisoning it. The result is exposed as an ordinary
\cpp{ProcessNoiseFunction} through
\cpp{ProcessNoise::GetProcessNoiseFunction}, so an adaptive filter differs
from a fixed one only in which callback was registered.

\section{Choosing a Filter}
\label{sec:flt-choice}

\Cref{tab:flt-comparison} summarizes the trade space. Two rules of thumb cover
most cases. First, if the state mixes quantities whose variances differ by more
than about eight orders of magnitude --- which any coupled orbit--clock state
does --- use a factorized filter; the dense EKF will silently lose the small
variances no matter how carefully the Joseph form is applied. Second, if the
measurement model spans two epochs, the choice is made for you: only
\cpp{UDUStochasticCloningEKF} represents the delayed-state dependency
correctly.

\begin{table}[htbp]
  \centering
  \small
  \caption{Comparison of the recursive filters. $n$ is the state dimension,
           $m$ the number of measurement components per epoch. Costs are the
           dominant term per epoch.}
  \label{tab:flt-comparison}
  \begin{tabularx}{\textwidth}{LLLL}
    \toprule
    Filter & Numerical robustness & Cost & When to use \\
    \midrule
    \cpp{EKF}
      & Moderate. Joseph form keeps $\mat{P}$ symmetric and PSD to first order,
        but the condition number of $\mat{P}$ is carried in full; small
        variances are lost to round-off.
      & $O(n^3)$
      & Default. Well-scaled states, moderate dynamic range, when the smoother
        and the diagnostics matter more than conditioning. \\
    \addlinespace
    \cpp{SchmidtEKF}
      & As \cpp{EKF}; the truncated gain is legal because the Joseph form
        tolerates a suboptimal $\mat{K}$.
      & $O(n^3)$
      & When a weakly observable bias must contribute its uncertainty without
        being solved for. \\
    \addlinespace
    \cpp{UKF}
      & Lowest. Short covariance form, no Joseph stabilization, no outlier
        gate; a Cholesky failure of \cref{eq:flt-ukf-chol} aborts the step.
      & $(2n{+}1)$ propagations, $O(n^3)$
      & Strong nonlinearity over the step, or a dynamics model with no usable
        Jacobian. Not for long unattended runs. \\
    \addlinespace
    \cpp{UDUEKF}
      & Highest among the covariance filters. PSD by construction, effective
        condition number $\sqrt{\kappa(\mat{P})}$, variances spanning many
        decades survive.
      & $O(n^3)$ time update, $O(n^2)$ per scalar measurement
      & Coupled orbit--clock estimation, long arcs, tight clock states.
        Requires diagonal $\mat{Q}$ and $\mat{R}$. \\
    \addlinespace
    \cpp{SRIF}
      & Highest overall. Backward-stable Householder updates, no gain matrix,
        admits an infinite prior variance (zero information).
      & $O(n^3)$ per step, plus an STM inverse
      & Cold starts with no meaningful a-priori state; full (non-diagonal)
        $\mat{R}$; batch-like arcs run sequentially. Needs an invertible
        $\mat{F}$. \\
    \addlinespace
    \cpp{UDUStochastic\-Cloning\-EKF}
      & As \cpp{UDUEKF}, on a $2n$ state.
      & $O((2n)^3)$ time update
      & Time-differenced carrier phase or any other delayed-state observable;
        required for a consistent TDCP smoother. \\
    \addlinespace
    \cpp{RunBatchFilter}
      & Good. LDLT normal equations with an SVD pseudo-inverse fallback for
        rank-deficient geometries.
      & $O(N n^2)$ per iteration
      & Snapshot positioning, filter initialization, and post-fit consistency
        checks. No dynamics, no process noise. \\
    \bottomrule
  \end{tabularx}
\end{table}

\section{Model Boundaries}
\label{sec:flt-boundaries}

\begin{itemize}[leftmargin=1.4em]
  \item \textbf{Diagonal noise requirements.} \cpp{UDUEKF} and
    \cpp{UDUStochasticCloningEKF} require a diagonal process noise $\mat{Q}$
    and a diagonal measurement noise $\mat{R}$: the Thornton time update
    weights the noise block by $\tilde{\mat{D}}$, and the Carlson update
    processes one scalar row at a time. Correlated process noise --- including
    the CWNA block of \cref{eq:flt-cwna} and the clock blocks of
    \cref{eq:flt-clock-q2,eq:flt-clock-q3} --- must be supplied through the
    mapping matrix $\mat{G}$. The time update asserts on a non-diagonal
    $\mat{Q}$; a non-diagonal $\mat{R}$ is \emph{not} asserted and its
    off-diagonal entries are silently ignored.

  \item \textbf{Linearization.} \cpp{EKF}, \cpp{SRIF}, \cpp{UDUEKF}, and the
    cloning filter are extended filters: they re-linearize about the current
    estimate at every step, and both $\mat{F}$ and $\mat{H}$ come from the
    registered callbacks. Their covariance is therefore correct only to first
    order in the state error, and a large initial error can produce a
    consistent-looking but wrong solution. The \cpp{UKF} removes the
    linearization but not the Gaussian assumption; none of the filters
    represents a multimodal posterior.

  \item \textbf{Gaussian, white, uncorrelated noise.} Every filter here assumes
    zero-mean Gaussian process and measurement noise, white in time and
    uncorrelated with each other. Coloured measurement noise (multipath,
    residual ionospheric delay, unmodelled solar radiation pressure) violates
    this and must be either estimated as a state, carried as a consider
    parameter (\cref{sec:flt-schmidt}), or absorbed by inflating $\mat{R}$.

  \item \textbf{Sequential-update ordering.} The Carlson recursion of
    \cref{sec:flt-udu} is exact for a linear model regardless of the order in
    which the scalar components are processed, but for a nonlinear model the
    result depends weakly on the ordering, since each row is linearized about
    the state produced by the preceding rows. The ordering rule of
    \cref{sec:flt-cloning} (delayed-state rows first) is a cost optimization,
    not a correctness requirement.

  \item \textbf{Outlier rejection.} The default gate is a per-component
    $\tau\sigma$ test against the filter's own innovation covariance; it has no
    protection against a slowly growing error that inflates $S_{ii}$ in step
    with the residual, and it screens components independently rather than
    searching for a consistent subset. It is shared by \cpp{EKF}, \cpp{SRIF},
    and \cpp{UDUEKF}; the \cpp{UKF} does not screen at all. Register a
    \cpp{FaultDetectionFunction} for anything stronger.

  \item \textbf{Smoothing.} \cpp{EKF::UpdateSmoother} implements the standard
    RTS recursion, which is exact only under the Markov assumption. With TDCP
    or any other delayed-state observable, the smoothed covariance it reports
    is optimistic; the correct gain is $\mat{J}_k$ of
    \cref{eq:flt-sc-gain}. The smoother also requires the full log
    (\cpp{InitializeLogger} must be called with the epoch count in advance) and
    is therefore memory-bound at $O(N n^2)$.

  \item \textbf{Batch estimator.} \cpp{RunBatchFilter} reports a covariance
    built from the measurement information only; the soft prior of
    \cref{eq:flt-batch-prior} regularizes the solution but is excluded from
    \cref{eq:flt-batch-cov}. A converged batch solution with a very weak
    geometry can therefore report a large covariance even though the estimate
    is pinned by the prior. Convergence is declared on the correction norm
    alone --- no line search, no trust region --- so a poor initial guess can
    diverge silently within \code{max\_iterations}.

  \item \textbf{Clock model fidelity.} \Cref{eq:flt-clock-q2,eq:flt-clock-q3}
    model the oscillator as two or three integrated Wiener processes. This
    captures the white-frequency, random-walk-frequency, and random-run
    branches of \cref{eq:flt-clock-allan} but \emph{not} flicker noise (the
    $\tau^0$ floor), deterministic frequency drift, or the temperature and
    radiation sensitivity of a real flight oscillator. The coefficients of
    \cref{tab:flt-clock} are nominal specification values, not measured unit
    data.
\end{itemize}
